\documentclass[11pt]{article}

\usepackage[margin=1in]{geometry}
\usepackage{amsmath, amssymb}
\usepackage{listings}
\usepackage{xcolor}
\usepackage{hyperref}
\usepackage{enumitem}
\usepackage{titlesec}

\definecolor{codebg}{RGB}{245,245,245}
\definecolor{codecomment}{RGB}{110,110,110}
\definecolor{codekeyword}{RGB}{0,0,150}

\lstset{
    backgroundcolor=\color{codebg},
    basicstyle=\ttfamily\small,
    commentstyle=\color{codecomment}\itshape,
    keywordstyle=\color{codekeyword}\bfseries,
    language=Python,
    breaklines=true,
    frame=single,
    tabsize=4,
    showstringspaces=false,
    xleftmargin=1em,
    xrightmargin=1em
}

\titleformat{\section}{\normalfont\Large\bfseries}{\thesection}{1em}{}
\titleformat{\subsection}{\normalfont\large\bfseries}{\thesubsection}{1em}{}

\newcommand{\complexity}[2]{\textbf{Time:} #1 \quad \textbf{Space:} #2}

\title{LeetCode Patterns Cheat Sheet}
\author{}
\date{}

\begin{document}
\maketitle
\tableofcontents

\section{Sliding Window}

Used for contiguous subarray/substring problems where a window expands and
contracts over the input while maintaining some invariant (sum, count,
distinct characters, etc.).

\subsection{Fixed-size window}
\complexity{$O(n)$}{$O(1)$ or $O(k)$}

\begin{lstlisting}
def max_sum_fixed_window(nums, k):
    window_sum = sum(nums[:k])
    best = window_sum
    for i in range(k, len(nums)):
        window_sum += nums[i] - nums[i - k]
        best = max(best, window_sum)
    return best
\end{lstlisting}

\subsection{Variable-size window (two pointers, same direction)}
\complexity{$O(n)$}{$O(1)$ or $O(\Sigma)$ for a character/count map}

\begin{lstlisting}
def smallest_subarray_with_sum_at_least(nums, target):
    left = 0
    total = 0
    best = float('inf')
    for right, val in enumerate(nums):
        total += val
        while total >= target:
            best = min(best, right - left + 1)
            total -= nums[left]
            left += 1
    return best if best != float('inf') else 0
\end{lstlisting}

\textbf{Recognize it when:} you need the longest/shortest/count of
contiguous subarrays or substrings satisfying a condition, and elements can
be added/removed incrementally.

\textbf{Classic problems:} Longest Substring Without Repeating Characters
(3), Minimum Window Substring (76), Longest Repeating Character
Replacement (424), Permutation in String (567), Max Consecutive Ones III
(1004), Fruit Into Baskets (904), Subarrays with K Different Integers
(992).

\section{Two Pointers}

Two indices moving through a (usually sorted) structure, either converging
from both ends or moving in the same direction at different speeds.

\subsection{Opposite ends (sorted array)}
\complexity{$O(n)$}{$O(1)$}

\begin{lstlisting}
def two_sum_sorted(nums, target):
    left, right = 0, len(nums) - 1
    while left < right:
        s = nums[left] + nums[right]
        if s == target:
            return [left, right]
        elif s < target:
            left += 1
        else:
            right -= 1
    return [-1, -1]
\end{lstlisting}

\textbf{Classic problems:} Two Sum II (167), 3Sum (15), Container With
Most Water (11), Trapping Rain Water (42), Valid Palindrome (125), Sort
Colors / Dutch National Flag (75).

\section{Fast \& Slow Pointers (Floyd's Cycle Detection)}

One pointer moves twice as fast as the other; useful for linked lists and
detecting cycles or finding midpoints.

\complexity{$O(n)$}{$O(1)$}

\begin{lstlisting}
def has_cycle(head):
    slow = fast = head
    while fast and fast.next:
        slow = slow.next
        fast = fast.next.next
        if slow is fast:
            return True
    return False
\end{lstlisting}

\textbf{Classic problems:} Linked List Cycle (141), Linked List Cycle II
(142), Middle of the Linked List (876), Happy Number (202), Find the
Duplicate Number (287).

\section{Prefix Sum \& Difference Arrays}

Precompute cumulative sums so range-sum queries become $O(1)$; a hash map
of prefix sums turns "subarray sum equals k" problems into lookups.

\complexity{$O(n)$ preprocessing, $O(1)$ per query}{$O(n)$}

\begin{lstlisting}
def subarray_sum_equals_k(nums, k):
    count = 0
    running = 0
    seen = {0: 1}
    for num in nums:
        running += num
        count += seen.get(running - k, 0)
        seen[running] = seen.get(running, 0) + 1
    return count
\end{lstlisting}

\textbf{Classic problems:} Subarray Sum Equals K (560), Range Sum Query -
Immutable (303), Product of Array Except Self (238), Continuous Subarray
Sum (523), Corporate Flight Bookings (1109, difference array).

\section{Binary Search}

Not just for sorted arrays -- also for "search on the answer" when the
answer space is monotonic (a feasibility check that flips from false to
true).

\subsection{Standard binary search}
\complexity{$O(\log n)$}{$O(1)$}

\begin{lstlisting}
def binary_search(nums, target):
    lo, hi = 0, len(nums) - 1
    while lo <= hi:
        mid = (lo + hi) // 2
        if nums[mid] == target:
            return mid
        elif nums[mid] < target:
            lo = mid + 1
        else:
            hi = mid - 1
    return -1
\end{lstlisting}

\subsection{Binary search on the answer}
\complexity{$O(n \log(\text{range}))$}{$O(1)$}

\begin{lstlisting}
def min_capacity(weights, days):
    def feasible(capacity):
        needed_days, load = 1, 0
        for w in weights:
            if load + w > capacity:
                needed_days += 1
                load = 0
            load += w
        return needed_days <= days

    lo, hi = max(weights), sum(weights)
    while lo < hi:
        mid = (lo + hi) // 2
        if feasible(mid):
            hi = mid
        else:
            lo = mid + 1
    return lo
\end{lstlisting}

\textbf{Recognize it when:} the problem asks for a minimum/maximum value
that satisfies some monotonic "is this good enough?" predicate.

\textbf{Classic problems:} Search in Rotated Sorted Array (33), Koko
Eating Bananas (875), Capacity To Ship Packages Within D Days (1011),
Split Array Largest Sum (410), Median of Two Sorted Arrays (4).

\section{BFS / DFS on Graphs and Trees}

Traversal patterns for exploring reachability, shortest paths (unweighted),
connected components, and tree structure.

\subsection{BFS (shortest path, level order)}
\complexity{$O(V + E)$}{$O(V)$}

\begin{lstlisting}
from collections import deque

def bfs_shortest_path(graph, start, target):
    queue = deque([(start, 0)])
    visited = {start}
    while queue:
        node, dist = queue.popleft()
        if node == target:
            return dist
        for nxt in graph[node]:
            if nxt not in visited:
                visited.add(nxt)
                queue.append((nxt, dist + 1))
    return -1
\end{lstlisting}

\subsection{DFS (connected components, backtracking traversal)}
\complexity{$O(V + E)$}{$O(V)$ recursion stack}

\begin{lstlisting}
def dfs(graph, node, visited):
    visited.add(node)
    for nxt in graph[node]:
        if nxt not in visited:
            dfs(graph, nxt, visited)
\end{lstlisting}

\textbf{Classic problems:} Number of Islands (200), Rotting Oranges (994),
Word Ladder (127), Clone Graph (133), Course Schedule (207, also
topological sort), Pacific Atlantic Water Flow (417).

\section{Topological Sort}

Ordering of nodes in a DAG so that every directed edge $u \to v$ has $u$
before $v$. Implemented via Kahn's algorithm (BFS with in-degrees) or DFS
with a post-order stack.

\complexity{$O(V + E)$}{$O(V + E)$}

\begin{lstlisting}
from collections import deque

def topo_sort(num_nodes, edges):
    graph = [[] for _ in range(num_nodes)]
    indegree = [0] * num_nodes
    for u, v in edges:
        graph[u].append(v)
        indegree[v] += 1

    queue = deque(i for i in range(num_nodes) if indegree[i] == 0)
    order = []
    while queue:
        node = queue.popleft()
        order.append(node)
        for nxt in graph[node]:
            indegree[nxt] -= 1
            if indegree[nxt] == 0:
                queue.append(nxt)
    return order if len(order) == num_nodes else []  # empty => cycle
\end{lstlisting}

\textbf{Classic problems:} Course Schedule (207), Course Schedule II
(210), Alien Dictionary (269), Minimum Height Trees (310).

\section{Backtracking}

Explore a decision tree, trying candidates and undoing (\emph{backtracking})
when a partial solution can't lead to a valid one. Used for
permutations/combinations, constraint satisfaction, and exhaustive search
with pruning.

\complexity{Exponential in general, pruning reduces it in practice}{$O(\text{depth})$}

\begin{lstlisting}
def subsets(nums):
    result = []
    path = []

    def backtrack(start):
        result.append(path[:])
        for i in range(start, len(nums)):
            path.append(nums[i])
            backtrack(i + 1)
            path.pop()

    backtrack(0)
    return result
\end{lstlisting}

\textbf{Classic problems:} Subsets (78), Permutations (46), Combination
Sum (39), N-Queens (51), Word Search (79), Palindrome Partitioning (131).

\section{Dynamic Programming}

Break a problem into overlapping subproblems with optimal substructure;
cache results (memoization or bottom-up table) to avoid recomputation.

\subsection{1D DP}
\complexity{$O(n)$}{$O(n)$, often reducible to $O(1)$}

\begin{lstlisting}
def house_robber(nums):
    prev, curr = 0, 0
    for num in nums:
        prev, curr = curr, max(curr, prev + num)
    return curr
\end{lstlisting}

\subsection{2D DP (grid / two sequences)}
\complexity{$O(mn)$}{$O(mn)$, often reducible to $O(n)$}

\begin{lstlisting}
def longest_common_subsequence(text1, text2):
    m, n = len(text1), len(text2)
    dp = [[0] * (n + 1) for _ in range(m + 1)]
    for i in range(1, m + 1):
        for j in range(1, n + 1):
            if text1[i - 1] == text2[j - 1]:
                dp[i][j] = dp[i - 1][j - 1] + 1
            else:
                dp[i][j] = max(dp[i - 1][j], dp[i][j - 1])
    return dp[m][n]
\end{lstlisting}

\subsection{Knapsack (0/1 and unbounded)}

\begin{lstlisting}
def knapsack_01(weights, values, capacity):
    dp = [0] * (capacity + 1)
    for w, v in zip(weights, values):
        for cap in range(capacity, w - 1, -1):  # reverse => each item once
            dp[cap] = max(dp[cap], dp[cap - w] + v)
    return dp[capacity]
\end{lstlisting}

\textbf{Recognize it when:} you're asked for a min/max/count over choices
with recursive structure, and brute force overlaps subproblems (Fibonacci
style, or "choose to include/exclude").

\textbf{Classic problems:} Climbing Stairs (70), House Robber (198),
Coin Change (322), Longest Increasing Subsequence (300), Edit Distance
(72), Partition Equal Subset Sum (416), Longest Common Subsequence (1143).

\section{Greedy}

Make the locally optimal choice at each step, provable to yield a globally
optimal result (usually via an exchange argument or matroid structure).
Often paired with sorting.

\complexity{$O(n \log n)$ (dominated by sort)}{$O(1)$ or $O(n)$}

\begin{lstlisting}
def can_attend_all_meetings(intervals):
    intervals.sort(key=lambda x: x[0])
    for i in range(1, len(intervals)):
        if intervals[i][0] < intervals[i - 1][1]:
            return False
    return True
\end{lstlisting}

\textbf{Classic problems:} Jump Game (55), Gas Station (134), Task
Scheduler (621), Non-overlapping Intervals (435), Minimum Number of
Arrows to Burst Balloons (452).

\section{Merge Intervals}

Sort intervals by start time, then merge or compare adjacent ones. A
specialization of the greedy pattern that shows up often enough to deserve
its own bucket.

\complexity{$O(n \log n)$}{$O(n)$}

\begin{lstlisting}
def merge_intervals(intervals):
    intervals.sort(key=lambda x: x[0])
    merged = [intervals[0]]
    for start, end in intervals[1:]:
        if start <= merged[-1][1]:
            merged[-1][1] = max(merged[-1][1], end)
        else:
            merged.append([start, end])
    return merged
\end{lstlisting}

\textbf{Classic problems:} Merge Intervals (56), Insert Interval (57),
Meeting Rooms II (253), Interval List Intersections (986).

\section{Monotonic Stack}

Maintain a stack that is always increasing or decreasing; used to find the
next/previous greater or smaller element in $O(n)$ instead of $O(n^2)$.

\complexity{$O(n)$}{$O(n)$}

\begin{lstlisting}
def daily_temperatures(temps):
    answer = [0] * len(temps)
    stack = []  # indices, temps decreasing
    for i, t in enumerate(temps):
        while stack and temps[stack[-1]] < t:
            j = stack.pop()
            answer[j] = i - j
        stack.append(i)
    return answer
\end{lstlisting}

\textbf{Classic problems:} Daily Temperatures (739), Next Greater Element
I/II (496, 503), Largest Rectangle in Histogram (84), Trapping Rain Water
(42, stack variant), Remove K Digits (402).

\section{Monotonic Deque (Sliding Window Maximum)}

Combine sliding window with a monotonic deque to get window max/min in
amortized $O(1)$ per step.

\complexity{$O(n)$}{$O(k)$}

\begin{lstlisting}
from collections import deque

def sliding_window_maximum(nums, k):
    dq = deque()  # indices, values decreasing
    result = []
    for i, num in enumerate(nums):
        while dq and nums[dq[-1]] < num:
            dq.pop()
        dq.append(i)
        if dq[0] <= i - k:
            dq.popleft()
        if i >= k - 1:
            result.append(nums[dq[0]])
    return result
\end{lstlisting}

\textbf{Classic problems:} Sliding Window Maximum (239), Shortest
Subarray with Sum at Least K (862).

\section{Heap / Top-K Elements}

A (min/max) heap keeps the $k$ best or worst elements seen so far without
sorting everything, and merges multiple sorted structures efficiently.

\complexity{$O(n \log k)$}{$O(k)$}

\begin{lstlisting}
import heapq

def top_k_frequent(nums, k):
    counts = {}
    for n in nums:
        counts[n] = counts.get(n, 0) + 1
    heap = []
    for num, freq in counts.items():
        heapq.heappush(heap, (freq, num))
        if len(heap) > k:
            heapq.heappop(heap)
    return [num for freq, num in heap]
\end{lstlisting}

\textbf{Classic problems:} Kth Largest Element in an Array (215), Top K
Frequent Elements (347), Merge k Sorted Lists (23), Find Median from Data
Stream (295), Task Scheduler (621).

\section{Union-Find (Disjoint Set)}

Track connected components under union operations, with path compression
and union by rank/size for near-$O(1)$ amortized operations.

\complexity{$O(\alpha(n))$ per operation}{$O(n)$}

\begin{lstlisting}
class UnionFind:
    def __init__(self, n):
        self.parent = list(range(n))
        self.rank = [0] * n

    def find(self, x):
        if self.parent[x] != x:
            self.parent[x] = self.find(self.parent[x])
        return self.parent[x]

    def union(self, a, b):
        ra, rb = self.find(a), self.find(b)
        if ra == rb:
            return False
        if self.rank[ra] < self.rank[rb]:
            ra, rb = rb, ra
        self.parent[rb] = ra
        if self.rank[ra] == self.rank[rb]:
            self.rank[ra] += 1
        return True
\end{lstlisting}

\textbf{Classic problems:} Number of Connected Components (323), Redundant
Connection (684), Accounts Merge (721), Number of Islands II (305).

\section{Trie (Prefix Tree)}

Tree structure for storing strings by shared prefixes, enabling fast
prefix lookups and autocomplete-style queries.

\complexity{$O(L)$ per insert/search, $L$ = word length}{$O(\Sigma \cdot N \cdot L)$}

\begin{lstlisting}
class TrieNode:
    def __init__(self):
        self.children = {}
        self.is_word = False

class Trie:
    def __init__(self):
        self.root = TrieNode()

    def insert(self, word):
        node = self.root
        for ch in word:
            node = node.children.setdefault(ch, TrieNode())
        node.is_word = True

    def search(self, word):
        node = self.root
        for ch in word:
            if ch not in node.children:
                return False
            node = node.children[ch]
        return node.is_word
\end{lstlisting}

\textbf{Classic problems:} Implement Trie (208), Word Search II (212),
Design Add and Search Words Data Structure (211), Replace Words (648).

\section{Bit Manipulation}

XOR/AND/OR tricks for finding unique elements, counting bits, and
subset/state compression (bitmask DP).

\begin{lstlisting}
def single_number(nums):
    result = 0
    for n in nums:
        result ^= n
    return result

def count_bits(n):
    return bin(n).count('1')
\end{lstlisting}

\textbf{Classic problems:} Single Number (136), Counting Bits (338),
Number of 1 Bits (191), Sum of Two Integers (371), Subsets via bitmask
(78 revisited).

\section{Choosing a Pattern -- Quick Reference}

\begin{itemize}[leftmargin=1.5em]
    \item \textbf{Contiguous subarray/substring, expand \& contract:}
        Sliding Window.
    \item \textbf{Sorted array, pair/triplet sums:} Two Pointers.
    \item \textbf{Linked list cycle or midpoint:} Fast \& Slow Pointers.
    \item \textbf{Range sum queries, "subarray sums to k":} Prefix Sum.
    \item \textbf{Sorted search space, or "minimize/maximize feasible
        value":} Binary Search (possibly on the answer).
    \item \textbf{Grid/graph reachability, shortest path unweighted:}
        BFS/DFS.
    \item \textbf{Ordering with dependencies:} Topological Sort.
    \item \textbf{All combinations/permutations, constraint search:}
        Backtracking.
    \item \textbf{Optimal value with overlapping subproblems:} Dynamic
        Programming.
    \item \textbf{Sort then make locally optimal choice:} Greedy /
        Merge Intervals.
    \item \textbf{Next greater/smaller element:} Monotonic Stack.
    \item \textbf{Window max/min:} Monotonic Deque.
    \item \textbf{Kth largest/smallest, frequency ranking:} Heap.
    \item \textbf{Group connectivity, cycle detection in undirected
        graph:} Union-Find.
    \item \textbf{Prefix-based string lookups:} Trie.
    \item \textbf{Uniqueness, subset enumeration, small state space:}
        Bit Manipulation.
\end{itemize}

\end{document}
